\section{One nonlinear return on a whole parameter interval}
\label{sec:selected-nonlinear-offering}
The local estimates must supply a return which varies continuously with
its input. Selecting a new support dictionary independently at every
parameter would not suffice. We explain how to obtain the exact integral
identity, both activity bounds and continuity in one dictionary on an
arbitrary compact parameter interval. The scale purchases used below are
fixed before that interval and the finite volume are chosen.

Fix a finite cutoff $N$, adjacent levels $m+1,m$ with $m<N$, and a
period. Write $p_f,p_c$ for the two activity budgets, $x_0$ for the
reference coarse coupling and $L$ for the block side. The fine reference
coupling is $\beta_F(x_0)=x_0+s_G\log L$. Fix the already selected
incoming and outgoing paths, their block bridge and the scalar shift
$\gamma=\alpha_{N,m}$. Let $I=[a,b]$ with $a\le b$, and let
$Q_x$, $x\in I$, be a continuous family of fine representations in
the permitted fine band. Continuity here includes the regular and
forced coefficient arrays in their fixed labelled norms. Each $Q_x$
has the selected fine path and satisfies the same history constraints.

We use the following purchased estimates on these data. The regular
principal price is at most $E/16$, the complementary regular price at
most $2E/16$, and the preliminary forced price at most $K(x)/8$, where
\[
 E=c_{\rm st}\delta_E(x_0),\qquad
 K(x)=K_{\rm forecast}(Q_x.\beta),\qquad
 \delta_E(x_0)=x_0^{-11/10},\quad c_{\rm st}>0.
\]
These are the principal and complementary bounds of the preceding local
expansion, in the actual coarse budget. The selected path has its full
quadratic, cubic, quartic and loop bounds, with width $\sigma_c>0$;
its loop allowance is the selected allowance at level $m$. The initial
action is admissible in this same budget. If $A_c$ denotes the
support exponent used before the finite symmetry average, the scale
purchases also give
\begin{equation}\label{eq:selected-offer-width}
 A_c=b_c+3\le a_{\rm earn},\qquad A_c/B_c\le\sigma_c/4.
\end{equation}
Here $b_c$ is the coarse support exponent, $B_c$ its radius and
$a_{\rm earn}$ the earned exponent of the local expansion. In particular,
these width and exponent conditions concern the same output path as the
bounds, rather than a less expensive substitute.

\subsection{Keeping the physical and auxiliary data fixed}
At each $x\in I$, apply the local expansion to $Q_x$ on the physical
torus. Its hypotheses are supplied as follows. The block equation and
both adjacent kernel bounds come from the selected horizon. Its
all-period positive lower bounds apply to the two selected kernel levels;
the bridge intertwining identities and the periodized rows are those
of that same horizon. The input representation supplies the small-field
and forced budgets and the fine band condition. The expansion is applied
with the selected derivative orders, spatial weights, earned exponent,
and its two error reserves. Its differentiated remainder bounds are
retained with its output, as well as the undifferentiated integral
identity. They cannot be recovered from that identity alone.

Use the expansion at $x=a$ as the reference dictionary. For every other
$x$, identify elementary labels by the identity: the forests, geometric
supports, sector labels and component types are unchanged. The path
identity and the fixed regular and forced input dictionaries make this
an admissible identification. Thus all output arrays can be read in the
reference dictionary. This is precisely the comparison used in
Section~\ref{sec:finite-volume-parameter-continuity}; no claim that an
arbitrary choice of coordinates depends continuously on $x$ is needed.

The auxiliary tori used for local comparisons carry the same block
bridge and the same two paths. On such a torus the reference fine
representation is the bare path action with the pale activity array:
its density is
\[
 {\bf1}_{\{\text{all fine sites healthy}\}}
 \exp\{-\beta\mathcal K_f-\mathcal L_f\}.
\]
Here $\mathcal K_f$ and $\mathcal L_f$ are the fine kinetic and loop
terms. The pale array has zero regular cost and its stated forced cost.
The fine path bounds and the clock inequality for the loop allowance
therefore give the required incoming budgets. The fine threshold,
radius and band are exactly those at $\beta_F(x_0)$. The all-period
kernel bounds apply on every such torus, so the local expansion and its
differentiated remainder bounds are available there too. The auxiliary
expansions provide local comparisons; the density reconstruction below
is performed on the physical torus.

\subsection{The finite short density and its complement}
The regular dictionary contains two kinds of entries: connected
interaction clusters and path atoms. Divide them into a retained
infinite part and an exhaustible short part. The selected interaction clusters have at most two polymer vertices,
each formed from at most two elementary path letters of word length
at most $B_c$. Their connected coefficient is nonzero and each
polymer is admissible. These letters have empty input roster: they
contain neither incoming regular insertions nor incoming forced
slots. Every cluster outside this particular class is retained in
the complementary array. Its regular price is at most $3E/16$ by
the principal and complementary estimates. The retained path part
consists of atoms of word length greater than $B_c$. Let
$\mathcal P_d$ contain the selected interaction clusters and the
short path atoms whose elementary indices lie below the finite
cutoff $d$. The panels are increasing and exhaust every nonretained entry. Leave
the forced array unchanged at this intermediate operation, and set the
regular entries of $\mathcal P_d$ to zero. Call the resulting array
$T_{x,d}$.

This complement is the sum of the retained infinite part and the
unselected short tail, entry by entry. For a selected short entry both
sides vanish; for a retained entry only the retained summand is present;
for every other entry only the tail is present. Thus the triangle
inequality adds their regular budgets without changing any support.
The path tail requires a separate estimate. If its original width is
$\sigma_c$ and its word length exceeds $B_c$, exponential summability
and $A_c/B_c\le\sigma_c/4$ give a tail bound
\[
 e^{-\sigma_cB_c/4}\mathcal C_{\rm path}\le g,
                  \qquad g=\beta_F(x_0)^{-23/20}.
\]
Here $\mathcal C_{\rm path}$ is the finite path-jet price at the chosen
stage. The first inequality follows by splitting the exponential row
weight into the part paying the support weight and the part paying
the tail; the remaining positive weight is discarded. The second is
one of the scale purchases, imposed after the stage is fixed. The
finite symmetry average costs $C_{\rm turn}$, so the retained total
has budget $3E/16+C_{\rm turn}g$.

The unselected short tail has budget at most $g$ for all sufficiently
large $d$, uniformly in $x\in I$. Here is the uniformity argument in
the actual jet norms. On the fixed torus changing an anchor costs at
most $((1+D)^p)^j$, where $D$ is its diameter and $p,j$ are the
spatial power and derivative order. The sum of the rooted regular
majorants is finite, since there are finitely many anchors. Thus at
any fixed $x$ there is a summable majorant in the reference anchors,
and exhaustion of the short labels makes its unselected sum tend to
zero. Norm continuity gives a neighborhood of $x$ where the difference
array has arbitrarily small total price, so the same truncation works
there. A finite subcover of $I$ gives one truncation for the whole
interval. Consequently
\begin{equation}\label{eq:selected-offer-truncation}
 \|T_{x,d}^{\rm reg}\|_{p_c}
       \le h_0:=\frac3{16}E+C_{\rm turn}g+g.
\end{equation}
The dependence of the tail index on the fixed torus and interval is
allowed; its bound $h_0$ uses the purchased, volume-independent budget.

Allocate each entry $F_a$ of the finite panel among its allowed
anchors with nonnegative weights $\omega_{oa}$ summing to one. For
an interaction cluster this weight is the reciprocal of the number
of its allowed anchors; for a path atom it is the corresponding
fixed local allocation. Average over signed lattice symmetries and
translations. Define the unmasked short density
\[
                   f_{x,d}(o,U)=\sum_{a\in\mathcal P_d}
                                      \omega_{oa}F_{x,a}(U).
\]
This is the finite connected-cluster contribution plus the finite
path-atom contribution, with the same scalar shift $\gamma$.
The allocation is local: a nonzero weight implies that the entry's
support lies in a fixed radius-$r$ neighborhood of $o$. The weights
sum to one even after the finite average. The masked panel contribution
is therefore the sum of
\[
 E_{x,d}(o,U)=\sum_{a\in\mathcal P_d}\omega_{oa}
             {\bf1}_{\{U\text{ healthy on }\operatorname{supp}F_a\}}
                                                       F_{x,a}(U).
\]

\subsection{Compatibility and the five auxiliary offers}
Use the auxiliary bare inputs described above on the five tori of
sides $2(w+i)$, $0\le i\le4$, with $w\ge2r$. Their coupling is
$Q_x.\beta$, and their path, budget, bridge, scalar shift and finite
label cutoff $d$ are the physical ones. Let $f^{(i)}_{x,d}$ denote
the corresponding short density. These are the actual auxiliary
offers to the five-period construction of
Section~\ref{sec:five-period-localization}.

We check its compatibility condition. Under replication to a larger
period, each short cluster meeting the designated anchor has a unique
lift within the fixed local patch. Conversely projection gives its
unique short cluster on the smaller period. The geometric margin
$4r+16\le |\Lambda|^{1/4}$ prevents wrapping in these identifications.
They preserve the elementary component types and the cutoff on their
labels, so they give a bijection of the finite panels at that anchor.
The local connected coefficient is preserved as follows. On globally
healthy fields the mask in each selected cluster equals one, so
its paid masked response equals the actual connected coefficient.
The graph coefficient depends only on the overlaps and equality
relations of the selected polymers; the local bijection preserves
both. Its factorial normalization is unchanged. The remaining
coefficient is the product of the corresponding polymer integrals.

For each polymer choose matching forests on the two periods. The
selected letters have empty roster, and the fields are globally
healthy, so the local integrals have no imposed input insertions or
flags. Let $I$ be the finite set of fluctuation coordinates read by
the larger-period integrand, and let $e:I\hookrightarrow J$ be its
embedding into the smaller-period coordinate set. The period margin
makes this map injective. The local windows correspond bijectively:
their principal matrices have identical entries, since the local
fields and paths agree. Their inverses therefore correspond, as do
their extensions by zero outside the windows. The covariance is the
symmetrized local inverse. Consequently, if $C'$ and $C$ are the two
covariances, then
\[
 C'|_{I\times I}=C[e(I),e(I)].
\]
This assertion concerns finite covariance submatrices; the full
Gaussian spaces can have different dimensions.

For completeness, a centered Gaussian with covariance $C$ has
characteristic function
\[
 \exp(-\langle t,Ct\rangle/2).
\]
Its projection to the coordinates $e(I)$ therefore has characteristic
function
\[
 \exp(-\langle t,C[e(I),e(I)]t\rangle/2).
\]
The covariance identity
shows that this projected law is the same as the larger Gaussian's
marginal on $I$. Locality expresses both integrands as the same
measurable function of these coordinates: the block tags, frame
tags, elementary path words, and transported mean agree. Thus their
Gaussian integrals agree. This equality holds for every pair of
outer Haar-slice values, because the selected, unrostered terms are
independent of those slice choices after the matching local data
are fixed. The inner integral is consequently a common constant;
integrating against the normalized outer Haar measures preserves
it. This proves equality of the local integrals, including their
normalizations.

The forests used by the expansion are averaged, so equality for one
matching pair must be lifted to equality of those averages. Forests
are generated from finite independent local choices. Restrict these
choices to the embedded smaller patch and project each local choice
to its smaller-period version. The resulting map $P$ between the
finite choice sets has constant-size fibers: restriction forgets
independent coordinates, and the local projection has constant-size
fibers at each retained coordinate. Hence for any function $F$ on
the smaller choice set $\mathcal R$,
\[
 \frac1{|\mathcal R'|}\sum_{r'\in\mathcal R'}F(P r')
                  =\frac1{|\mathcal R|}\sum_{r\in\mathcal R}F(r).
\]
Locality lets one replace an arbitrary larger forest by the matching
lift of its restriction without changing the polymer integral.
Apply the displayed identity first to the block and frame tags,
then to the forest choices. This proves equality of the actual
averaged polymer coefficients. The normalized anchor weight is also
preserved because its allowed-anchor set is in bijection. For path atoms the identity follows directly from
replicating the same word of link increments. Finally the finite
symmetry average commutes with replication: translation fibers all
have equal cardinality, which cancels against the normalized average.
Thus the short density is preserved on globally healthy replicated
fields. Passing two tori to their product period and using locality
there proves the required compatibility for equal local patterns.
The auxiliary bare input and the physical input can be compared here
because the short cluster construction uses the same path and the
same healthy local generation; terms retained from the incoming
regular array remain in the complementary part, rather than being
identified with bare-input terms.

To price the centered auxiliary total, use its exact decomposition
into the compensated interacting action, the complementary regular
array and a constant. The first term has the healthy derivative
bound of Section~\ref{sec:healthy-boundary-jets}. For an auxiliary
torus $\Lambda_i$, its complementary array has budget
$C_{\rm turn}(g+g)+g$ after the same uniform-tail choice. Summing over
its anchors bounds every allowed derivative of the total by
$|\Lambda_i|[C_{\rm turn}(g+g)+g]$. Centering subtracts its flat value
and costs at most a further copy of this bound. Consequently the
centered short total has derivative price
\[
 q_i=\mathcal T_j(g)
             +2|\Lambda_i|[C_{\rm turn}(g+g)+g],
\]
where $\mathcal T_j(g)$ is the healthy-boundary derivative price of
order $j$ for the compensated action. The finite-volume continuity
argument applies to the entire closed fine coupling band on each of
these five fixed tori, so one $d$ suffices for all their couplings.
The comparison charts are admissible by the purchased local gauge
extension and logarithmic-radius margin. These are precisely the
hypotheses used in the boundary-jet estimate, and apply to the same
compensated action and selected stationary path.

Let $g_{x,d}(o,U)$ be the folded counterterm made from these five
centered offers. Formula~\eqref{eq:five-period-local-price} gives its
jet bound $c'=\mathcal O_z(q_0,\ldots,q_4)$ on the radius-$z$
healthy neighborhood. The paid gauge extension puts every such
healthy folded field into the chart just used; gauge invariance
preserves the bound. Let $M$ bound the coarse covering number of
this neighborhood. Its weighted regular cost per anchor is
\[
                              s'=c'e^{A_cM}.
\]
This follows by multiplying its jet bound by its support weight,
which is at most $e^{A_cM}$. It introduces no volume factor.

\subsection{The discrepancy is a forced term}
Use the shifted folded sum $h_{x,d}(o,U)$ of
\eqref{eq:five-period-shifted-localization}, and define
\[
 D_{x,d}(o,U)=E_{x,d}(o,U)-f_{x,d}(o,1)-h_{x,d}(o,U).
\]
The exact identity \eqref{eq:five-period-exact-decomposition} becomes
\begin{equation}\label{eq:selected-offer-exact-split}
 A_x(U)=A_{T_{x,d}}(U)+|\Lambda|f_{x,d}(0,1)
       +\sum_o\chi_{N_z(o)}(U)g_{x,d}(o,U)
       +\sum_oD_{x,d}(o,U).
\end{equation}
Here $A_x$ is the original regular sum, and $\chi_{N_z(o)}$ is its
healthy indicator. The formula holds at every field. On a healthy
radius-$2z$ neighborhood the allocation supports are healthy, and
the folded identity identifies $h$ with the centered $f$. Hence
$D_{x,d}(o,U)=0$ there. The discrepancy is local on that enlarged
neighborhood and has the gauge invariance and measurability inherited
from its finitely many summands. Its value bound is obtained by the
triangle inequality for the panel sum and the five folded sums; this
is the local value estimate purchased along with the original
expansion. Write $q$ for this value bound, and $M_2$ for the covering-number
bound of the doubled support. Its soft absorption cost is
\[
 b_{\rm soft}=\operatorname{ballCost}(2z)\,q
                             e^{(b_c+2)M_2}.
\]
The factor $\operatorname{ballCost}(2z)$ bounds the number of anchors
whose discrepancy support can meet a fixed site; the extra exponent
reserve is the one needed for the raised forced budget. Finite
symmetry averaging preserves the corresponding normalized anchor
sum. The local value estimates and logarithmic-radius geometry give
$b_{\rm soft}\le\beta_F(x_0)^{-1/8}$ after increasing the entry
threshold. This is a separate purchase from the regular jet cost
$s'=c'e^{A_cM}$: $q$ bounds the discrepancy on arbitrary fields,
whereas $c'$ bounds the folded counterterm's healthy jets.

Crucially, the last sum in \eqref{eq:selected-offer-exact-split} is
not discarded. Put $-D_{x,d}$ into the logarithmic correction array,
and absorb it into the forced gas by the exact identity
$e^{-D}=1+(e^{-D}-1)$ and the support-join construction of
Section~\ref{sec:soft-log-fusion}. The vanishing just proved supplies
the eligibility condition of each new forced term. The old forced
gas and this correction are combined with their actual support and
inventory labels. Their hard input budget is at most $K(x)/8$,
and their soft absorption cost is at most $\beta_F(x_0)^{-1/8}$;
the purchased absorption inequalities give forced output budget
$K(x)$. This operation changes the forced gas exactly by the factor
$\exp(-\sum_oD_{x,d}(o,U))$.

The regular output consists of $T_{x,d}$ and the folded counterterms,
followed by the fixed symmetry average. Its budget is at most
$h_0+C_{\rm turn}s'$. The scale purchase gives
\[
 C_{\rm turn}g+g+C_{\rm turn}c'e^{A_cM}\le E/2.
\]
To justify that this purchase is available, $g$ and $c'$ have upper
exponent $-23/20$, while $E$ has exponent $-11/10$. Their fixed
stage constants and powers of the logarithmic radius do not change
this strict inequality of exponents. Thus an entry threshold makes
the displayed inequality hold for every later coupling. Combining
it with \eqref{eq:selected-offer-truncation} gives the bound
$11E/16\le E$. The forced estimate above is retained separately.

\subsection{The exact return and its parameter dependence}
Choose one spanning tree on the physical torus, the canonical sector
measure, and the bookings associated with the selected expansion.
Keep them fixed on $I$. Let $\mathcal X$ be the enlarged reference
dictionary, and let $v_x$ be the array formed from the truncated local
expansion, the local counterterms and the forced correction just
constructed. The scalar correction is
\[
 \tau_x=\text{the original scalar payment at }x
             -\text{the constant term removed by the local subtraction}.
\]
This sign follows by moving the removed constant from the action to
its vacuum parameter. The exact decomposition of the local expansion
and the local subtraction identity now give, for every coarse
configuration, the original Gaussian return evaluated using
$(\mathcal X,v_x,\tau_x)$. This is equality of the complete
unnormalized density, with the selected $\gamma$, not merely equality
of a finite jet at the identity. Absolute summability of the regular
arrays justifies subtraction and regrouping: the known and subtracted
regular occurrences add term by term to the original occurrences, and
the discrepancy factor is included in the forced gas as in
\eqref{eq:selected-offer-exact-split}. The finite regrouping and
forced absorption preserve the integral identity by their exact
exponential and gas identities.

All regular entries are local and gauge invariant with the prescribed
jets; the forced entries are local, measurable and gauge invariant.
Their bounds are $E$ and $K(x)$ respectively. The common-dictionary
parameter comparison survives truncation, addition, fixed pullbacks and
finite averages, each of which is continuous in these array norms.
For the remaining infinite arrays, the summable comparison majorants
justify passage to the limit. Thus $v_x$ has the required continuous
parameter dependence in the same dictionary $\mathcal X$.

We have proved the interval return needed by the scalar coupling reset:
for every compact interval and every allowed input curve there is one
dictionary, one tree and one sector choice, with an exact return and
both local budgets throughout the interval. This statement still uses
the simultaneous scale purchases listed at the start. Their selection
must precede all intervals, histories and periods. The scalar reset
then has a separate task: choose the new coupling and vacuum parameters
without losing this density identity or its labelled estimates.
